% D.31 integrated after checkpoint09 under the parent's authorization.
\subsection{A unit criterion for the remaining finite cusps}

\begin{claim}{D.31}{Proved}\label{clm:D31}
Let $X/\mathbf Q$ be a smooth geometrically connected
hyperbolic curve with smooth completion $C$, and let $s$
be a continuous noncuspidal section. Suppose its finite-place
localizations are ordinary or cuspidal, and are ordinary
integral outside a finite set.
Let $T\subset C\setminus X$ be the set of closed boundary
points that occur as centres of local cuspidal sections.
If some nonconstant $u\in\mathcal O(X)^\times$ is regular
and nonzero at every point of $T$, then $s$ is a rational
point section.

More generally, let $r$ be the number of finite places
where $s$ is cuspidal. If some connected finite \'etale
neighbourhood $Y/\mathbf Q$ of $s$ satisfies
\[
          \operatorname{rank}
          \bigl(\mathcal O(Y)^\times/\mathbf Q^\times\bigr)>r,
\]
then $s$ is a rational point section. Existence of such a
neighbourhood is an additional hypothesis.

For the tripod
$U=\mathbf P^1_{\mathbf Q}\setminus\{0,1,\infty\}$,
cofinite ordinary integrality and the occurrence of at
most one cusp label already imply that $s$ is a rational
point section. In particular this conclusion applies
under D.30's finite-ramification hypothesis if its
exceptional local cusps have at most one label.
\end{claim}

\begin{proof}
Pull back the Kummer class of $u$:
\[
                A=s^\ast\kappa(u)
                   \in\widehat{\mathbf Q^\times}.
\]
At an ordinary place $p$, its localization is the class
of $u(x_p)\in\mathbf Q_p^\times$.
At a cuspidal place, $u$ extends as a regular nonzero
function at the centre. Its Kummer covers are then
unramified along the horizontal cusp, so the cusp inertia
acts trivially and every continuous splitting gives the
actual leading value $u(c_p)\in\mathbf Q_p^\times$.
This assertion includes arbitrary completed tangential
classes; it does not require a rational tangent vector.

After excluding finitely many model primes, $u$ extends
as a unit and all ordinary integral local values are
units. D.3 therefore gives $A=\kappa(a)$ for
$a\in\mathbf Q^\times$.
Injectivity of local Kummer theory shows
$u(x_p)=a$ at each ordinary place and $u(c_p)=a$
at each cuspidal place.
The nonconstant function $u$ extends to a finite map
$C\to\mathbf P^1$. All finite-place support points
belong to the finite closed fibre $u^{-1}(a)$.
The finite-place finite-support theorem used in D.5 and
D.22, namely Stix's Corollary~5 with its audited
finite-descent input, makes $s$ a point or cuspidal
section \citep{D_Stix2015}. Noncuspidality places the
point in $X(\mathbf Q)$.

For the neighbourhood assertion, its lifted section is
still noncuspidal, and ordinary integrality transfers
outside finitely many model primes by the finite
normalization comparison in D.25--D.26.
An ordinary local point lifts to an ordinary point;
a boundary point lifts only to boundary points.
Thus there are at most $r$ relevant closed boundary
centres on $Y$.
The unit lattice embeds by its divisor in the free
abelian group on the boundary. The map recording its
orders at those at most $r$ centres has nonzero kernel
when its rank exceeds $r$. A nonzero lattice element
in that kernel is a nonconstant unit regular and
nonzero at all relevant centres. The first assertion
applied to $Y$ proves that the lifted section, hence
$s$, is geometric.

For the tripod, if a cusp label $c$ occurs, use
\[
\begin{array}{c|ccc}
c&0&1&\infty\\ \hline
u(t)&1-t&t&t/(t-1)\\
u(c)&1&1&1 .
\end{array}
\]
If it is the only label, each selected unit satisfies
the first assertion's hypothesis. There is also a
direct contradiction if any cusp occurs: D.3 makes
its global class rational, localization at that cusp
makes the rational value equal to one, whereas none
of the selected functions takes the value one on $U$.
Ordinary places exist by cofinite integrality.
Hence no cuspidal places remain, and D.5 applies.
\end{proof}

\proofparagraph{The surviving arithmetic condition for several labels.}
Choose a finite set $S$ containing the archimedean place,
all nonintegral ordinary places and all cuspidal places,
and put
\[
 \Gamma=\mathbf Z[S_f^{-1}]^\times,\qquad
 A=s^\ast\kappa(t),\quad B=s^\ast\kappa(1-t).
\]
Their valuation coordinates vanish outside $S$, so
D.3 identifies $A,B$ with elements of $\widehat\Gamma$.
For $p\notin S$ their localizations are actual units
$(x_p,1-x_p)$.
Since $\widehat\Gamma^2$ is compact, its image in
$\prod_{p\notin S}(\mathbf Z_p^\times)^2$ is precisely
the closure of $\Gamma^2$. Thus the induced tuple lies in
\[
 \overline{\Gamma^2}\ \cap\
 \prod_{p\notin S}
 \{(x,y)\in(\mathbf Z_p^\times)^2:x+y=1\}.
 \tag{D.31.1}
\]
The extra cusp relations are $B_p=1$ at label zero,
$A_p=1$ at label one, and $A_p/B_p=-1$ at infinity.
The respective potentially nonintegral profinite
valuation pairs are $(\alpha_p,0)$, $(0,\beta_p)$,
and $(\gamma_p,\gamma_p)$.
They have not been proved to be pairs of ordinary integers.

Harari--Voloch's Theorem~1 identifies the closure condition
in \textup{(D.31.1)} with its outside-$S$ finite abelian
descent condition; their Conjecture~2 asserts the required
diagonal integral-point conclusion
\citep[Theorem~1, Conjecture~2, Remark~2.5]{D_HarariVoloch2010}.
If that conclusion holds for this tuple, D.26's compatible
almost-everywhere Kummer detection gives
$A=\kappa(a)$ and $B=\kappa(1-a)$ globally.
Each possible cusp relation then gives a contradiction,
and D.5 finishes. No such general diagonal-intersection
assertion is proved here.

Helm--Voloch's Theorem~4 gives an integral elliptic
realization once a common finite ramification set and
the local good elliptic curves outside it are supplied.
Its modular-curve consequence, Theorem~3, asserts
existence of an integral point on a modular twist,
not equality with the given tuple of local points
\citep[Theorems~3--4]{D_HelmVoloch2011}.
Nor are arbitrary finite \'etale neighbourhoods required
to be congruence modular curves.
Liu--Xu's projected all-place Brauer--Manin theorem
requires an ordinary completion at the omitted places;
their stronger outside-$S$ conclusion over $\mathbf Q$
is stated only when $S$ is archimedean
\citep[Definition~1.1, Theorems~1.2--1.4, Corollary~5.11]
{D_LiuXu2015}.
These theorems do not eliminate the multiple-label case
left by D.30 and D.31.
